\documentclass[12pt]{article}
\usepackage[UTF8]{inputenc}
\usepackage{xeCJK}
\setCJKmainfont{Sarabun}
\usepackage{enumitem}
\usepackage{geometry}
\geometry{a4paper, margin=2.5cm}
\usepackage{setspace}
\onehalfspacing

\begin{document}

\begin{center}
{\Large\textbf{ฟิชชันและฟิวชัน (เพิ่มเติม)}}\\[6pt]
{\normalsize แบบฝึกหัดเพิ่มเติม}\\[6pt]
{\normalsize วิชาฟิสิกส์ ม.ปลาย บทที่ 29}\\[4pt]
{\footnotesize ทั้งหมด 30 ข้อ}
\end{center}
\hrule
\bigskip

\section*{แบบฝึกหัดเพิ่มเติม}
\begin{enumerate}[label=\arabic*., itemsep=0.5\baselineskip, topsep=0.5\baselineskip, leftmargin=*]
\item ผลิตภัณฑ์หลักของฟิชชัน ²³⁵U มักประกอบด้วยธาตุในกลุ่มใด?
\begin{enumerate}[label=\alph*)., itemsep=0.3\baselineskip, topsep=0.3\baselineskip]
\item ก\ 
\begin{minipage}[t]{0.9\linewidth}
\raggedright แก๊สเฉื่อย
\end{minipage}
\item ข\ 
\begin{minipage}[t]{0.9\linewidth}
\raggedright ธาตุหนักปานกลาง (เช่น Ba, Kr)
\end{minipage}
\item ค\ 
\begin{minipage}[t]{0.9\linewidth}
\raggedright ธาตุเบา (เช่น He, Li)
\end{minipage}
\item ง\ 
\begin{minipage}[t]{0.9\linewidth}
\raggedright โลหะแอลคาไล
\end{minipage}
\end{enumerate}
\vspace{-0.3\baselineskip}
\begin{small}
\textit{คำตอบ: ข}
\end{small}

\item แหล่งพลังงานหลักของดวงอาทิตย์มาจากปฏิกิริยาอะไร?
\begin{enumerate}[label=\alph*)., itemsep=0.3\baselineskip, topsep=0.3\baselineskip]
\item ก\ 
\begin{minipage}[t]{0.9\linewidth}
\raggedright ฟิชชันของยูเรเนียม
\end{minipage}
\item ข\ 
\begin{minipage}[t]{0.9\linewidth}
\raggedright ฟิวชันของไฮโดรเจนเป็นฮีเลียม
\end{minipage}
\item ค\ 
\begin{minipage}[t]{0.9\linewidth}
\raggedright การสลายตัวของคาร์บอน
\end{minipage}
\item ง\ 
\begin{minipage}[t]{0.9\linewidth}
\raggedright การเผาไหม้เคมี
\end{minipage}
\end{enumerate}
\vspace{-0.3\baselineskip}
\begin{small}
\textit{คำตอบ: ข}
\end{small}

\item ค่า k ในทฤษฎีปฏิกิริยาลูกโซ่คืออะไร?
\begin{enumerate}[label=\alph*)., itemsep=0.3\baselineskip, topsep=0.3\baselineskip]
\item ก\ 
\begin{minipage}[t]{0.9\linewidth}
\raggedright จำนวนนิวตรอนที่ปล่อยต่อฟิชชัน
\end{minipage}
\item ข\ 
\begin{minipage}[t]{0.9\linewidth}
\raggedright ค่าประสิทธิผลนิวตรอน (neutron multiplication factor)
\end{minipage}
\item ค\ 
\begin{minipage}[t]{0.9\linewidth}
\raggedright อัตราการสลายตัว
\end{minipage}
\item ง\ 
\begin{minipage}[t]{0.9\linewidth}
\raggedright มวลวิกฤตของสาร
\end{minipage}
\end{enumerate}
\vspace{-0.3\baselineskip}
\begin{small}
\textit{คำตอบ: ข}
\end{small}

\item สารชะลอนิวตรอน (Moderator) ทำงานอย่างไร?
\begin{enumerate}[label=\alph*)., itemsep=0.3\baselineskip, topsep=0.3\baselineskip]
\item ก\ 
\begin{minipage}[t]{0.9\linewidth}
\raggedright เพิ่มพลังงานนิวตรอน
\end{minipage}
\item ข\ 
\begin{minipage}[t]{0.9\linewidth}
\raggedright ลดพลังงานนิวตรอนผ่านการชนอะตอม
\end{minipage}
\item ค\ 
\begin{minipage}[t]{0.9\linewidth}
\raggedright ดูดซับนิวตรอน
\end{minipage}
\item ง\ 
\begin{minipage}[t]{0.9\linewidth}
\raggedright เปลี่ยนนิวตรอนเป็นโปรตอน
\end{minipage}
\end{enumerate}
\vspace{-0.3\baselineskip}
\begin{small}
\textit{คำตอบ: ข}
\end{small}

\item ฟิชชันของ ²³⁵U ปล่อยนิวตรอนกี่ตัวต่อฟิชชันโดยเฉลี่ย?
\begin{enumerate}[label=\alph*)., itemsep=0.3\baselineskip, topsep=0.3\baselineskip]
\item ก\ 
\begin{minipage}[t]{0.9\linewidth}
\raggedright 1 ตัว
\end{minipage}
\item ข\ 
\begin{minipage}[t]{0.9\linewidth}
\raggedright 2-3 ตัว
\end{minipage}
\item ค\ 
\begin{minipage}[t]{0.9\linewidth}
\raggedright 5-6 ตัว
\end{minipage}
\item ง\ 
\begin{minipage}[t]{0.9\linewidth}
\raggedright 10 ตัว
\end{minipage}
\end{enumerate}
\vspace{-0.3\baselineskip}
\begin{small}
\textit{คำตอบ: ข}
\end{small}

\item สารฟิชชันที่ใช้ในระเบิดปรมาณูคืออะไร?
\begin{enumerate}[label=\alph*)., itemsep=0.3\baselineskip, topsep=0.3\baselineskip]
\item ก\ 
\begin{minipage}[t]{0.9\linewidth}
\raggedright ยูเรเนียม-238
\end{minipage}
\item ข\ 
\begin{minipage}[t]{0.9\linewidth}
\raggedright ยูเรเนียม-235 หรือพลูโทเนียม-239 ที่มีความบริสุทธิ์สูง
\end{minipage}
\item ค\ 
\begin{minipage}[t]{0.9\linewidth}
\raggedright ทอเรียม-232
\end{minipage}
\item ง\ 
\begin{minipage}[t]{0.9\linewidth}
\raggedright แคดเมียม
\end{minipage}
\end{enumerate}
\vspace{-0.3\baselineskip}
\begin{small}
\textit{คำตอบ: ข}
\end{small}

\item ข้อจำกัดหลักของการทำฟิวชันบนโลกคืออะไร?
\begin{enumerate}[label=\alph*)., itemsep=0.3\baselineskip, topsep=0.3\baselineskip]
\item ก\ 
\begin{minipage}[t]{0.9\linewidth}
\raggedright ไม่มีเชื้อเพลิง
\end{minipage}
\item ข\ 
\begin{minipage}[t]{0.9\linewidth}
\raggedright ยากที่จะสร้างและรักษาอุณหภูมิหลายร้อยล้านเคลวินได้นานพอ
\end{minipage}
\item ค\ 
\begin{minipage}[t]{0.9\linewidth}
\raggedright ไม่มีหลักการทางฟิสิกส์รองรับ
\end{minipage}
\item ง\ 
\begin{minipage}[t]{0.9\linewidth}
\raggedright กฎหมายห้าม
\end{minipage}
\end{enumerate}
\vspace{-0.3\baselineskip}
\begin{small}
\textit{คำตอบ: ข}
\end{small}

\item ทำไมยูเรเนียม-238 ไม่สามารถรักษาปฏิกิริยาลูกโซ่ด้วยตัวเองได้?
\begin{enumerate}[label=\alph*)., itemsep=0.3\baselineskip, topsep=0.3\baselineskip]
\item ก\ 
\begin{minipage}[t]{0.9\linewidth}
\raggedright ไม่ดูดซับนิวตรอน
\end{minipage}
\item ข\ 
\begin{minipage}[t]{0.9\linewidth}
\raggedright นิวตรอนเร็วที่ปล่อยไม่สามารถทำให้ ²³⁸U ฟิชชันได้เพียงพอ
\end{minipage}
\item ค\ 
\begin{minipage}[t]{0.9\linewidth}
\raggedright ²³⁸U สลายตัวเร็วเกินไป
\end{minipage}
\item ง\ 
\begin{minipage}[t]{0.9\linewidth}
\raggedright ไม่มีนิวตรอน
\end{minipage}
\end{enumerate}
\vspace{-0.3\baselineskip}
\begin{small}
\textit{คำตอบ: ข}
\end{small}

\item Breeder reactor มีความสามารถอย่างไร?
\begin{enumerate}[label=\alph*)., itemsep=0.3\baselineskip, topsep=0.3\baselineskip]
\item ก\ 
\begin{minipage}[t]{0.9\linewidth}
\raggedright ผลิตไฟฟ้าจากฟิชชันเท่านั้น
\end{minipage}
\item ข\ 
\begin{minipage}[t]{0.9\linewidth}
\raggedright ผลิตเชื้อเพลิงใหม่จากสารไม่ fissile
\end{minipage}
\item ค\ 
\begin{minipage}[t]{0.9\linewidth}
\raggedright ชะลอนิวตรอนให้ช้าลงเท่านั้น
\end{minipage}
\item ง\ 
\begin{minipage}[t]{0.9\linewidth}
\raggedright ทำให้ขยะนิวเคลียร์หมดไป
\end{minipage}
\end{enumerate}
\vspace{-0.3\baselineskip}
\begin{small}
\textit{คำตอบ: ข}
\end{small}

\item ขยะนิวเคลียร์จากเครื่องปฏิกรณ์พลังงานนิวเคลียร์แบ่งระดับตามกัมมันตภาพรังสีได้กี่ระดับ?
\begin{enumerate}[label=\alph*)., itemsep=0.3\baselineskip, topsep=0.3\baselineskip]
\item ก\ 
\begin{minipage}[t]{0.9\linewidth}
\raggedright 2 ระดับ
\end{minipage}
\item ข\ 
\begin{minipage}[t]{0.9\linewidth}
\raggedright 3 ระดับ
\end{minipage}
\item ค\ 
\begin{minipage}[t]{0.9\linewidth}
\raggedright 4 ระดับ
\end{minipage}
\item ง\ 
\begin{minipage}[t]{0.9\linewidth}
\raggedright 5 ระดับ
\end{minipage}
\end{enumerate}
\vspace{-0.3\baselineskip}
\begin{small}
\textit{คำตอบ: ข}
\end{small}

\item ปฏิกิริยาฟิชชันของ ²³⁵U ปล่อยพลังงานเฉลี่ย 200 MeV ถ้ากำลังการผลิตไฟฟ้าของโรงไฟฟ้านิวเคลียร์ 1000 MW จงหาว่าต้องใช้ฟิชชันกี่ครั้งต่อวินาที?
\begin{enumerate}[label=\alph*)., itemsep=0.3\baselineskip, topsep=0.3\baselineskip]
\item ก\ 
\begin{minipage}[t]{0.9\linewidth}
\raggedright 3.125 × 10¹⁹ ครั้ง/s
\end{minipage}
\item ข\ 
\begin{minipage}[t]{0.9\linewidth}
\raggedright 3.125 × 10¹⁶ ครั้ง/s
\end{minipage}
\item ค\ 
\begin{minipage}[t]{0.9\linewidth}
\raggedright 3.125 × 10¹³ ครั้ง/s
\end{minipage}
\item ง\ 
\begin{minipage}[t]{0.9\linewidth}
\raggedright 3.125 × 10²⁰ ครั้ง/s
\end{minipage}
\end{enumerate}
\vspace{-0.3\baselineskip}
\begin{small}
\textit{คำตอบ: ข}
\end{small}

\item ความแตกต่างหลักระหว่างระเบิดปรมาณูและระเบิดไฮโดรเจนคืออะไร?
\begin{enumerate}[label=\alph*)., itemsep=0.3\baselineskip, topsep=0.3\baselineskip]
\item ก\ 
\begin{minipage}[t]{0.9\linewidth}
\raggedright ระเบิดปรมาณูใช้ฟิวชัน ระเบิดไฮโดรเจนใช้ฟิชชัน
\end{minipage}
\item ข\ 
\begin{minipage}[t]{0.9\linewidth}
\raggedright ระเบิดปรมาณูใช้ฟิชชัน ระเบิดไฮโดรเจนใช้ฟิวชันเป็นหลัก
\end{minipage}
\item ค\ 
\begin{minipage}[t]{0.9\linewidth}
\raggedright ทั้งสองใช้ฟิชชันเหมือนกัน
\end{minipage}
\item ง\ 
\begin{minipage}[t]{0.9\linewidth}
\raggedright ระเบิดไฮโดรเจนมีขนาดเล็กกว่า
\end{minipage}
\end{enumerate}
\vspace{-0.3\baselineskip}
\begin{small}
\textit{คำตอบ: ข}
\end{small}

\item เหตุการณ์ใดที่เครื่องปฏิกรณ์นิวเคลียร์ฟุกะชิมะเกิดขึ้น?
\begin{enumerate}[label=\alph*)., itemsep=0.3\baselineskip, topsep=0.3\baselineskip]
\item ก\ 
\begin{minipage}[t]{0.9\linewidth}
\raggedright ระเบิดปรมาณู
\end{minipage}
\item ข\ 
\begin{minipage}[t]{0.9\linewidth}
\raggedright Meltdown จากระบบทำความเย็นล้มเหลวหลังแผ่นดินไหวและคลื่นสึนามิ
\end{minipage}
\item ค\ 
\begin{minipage}[t]{0.9\linewidth}
\raggedright การรั่วไหลของขยะนิวเคลียร์
\end{minipage}
\item ง\ 
\begin{minipage}[t]{0.9\linewidth}
\raggedright การจุดชนวนผิดพลาด
\end{minipage}
\end{enumerate}
\vspace{-0.3\baselineskip}
\begin{small}
\textit{คำตอบ: ข}
\end{small}

\item ทำไมปฏิกิริยาฟิวชันจึงให้พลังงาน?
\begin{enumerate}[label=\alph*)., itemsep=0.3\baselineskip, topsep=0.3\baselineskip]
\item ก\ 
\begin{minipage}[t]{0.9\linewidth}
\raggedright เพราะผลรวมมวลของผลิตภัณฑ์น้อยกว่าผลรวมมวลของสารตั้งต้น
\end{minipage}
\item ข\ 
\begin{minipage}[t]{0.9\linewidth}
\raggedright เพราะผลรวมมวลของผลิตภัณฑ์มากกว่าผลรวมมวลของสารตั้งต้น
\end{minipage}
\item ค\ 
\begin{minipage}[t]{0.9\linewidth}
\raggedright เพราะเกิดปฏิกิริยาเคมีระหว่างนิวเคลียส
\end{minipage}
\item ง\ 
\begin{minipage}[t]{0.9\linewidth}
\raggedright เพราะนิวเคลียสสลายตัว
\end{minipage}
\end{enumerate}
\vspace{-0.3\baselineskip}
\begin{small}
\textit{คำตอบ: ก}
\end{small}

\item จงหาว่าถ้าต้องการให้ปฏิกิริยาลูกโซ่ดำเนินต่อเนื่อง (k=1) โดยมีนิวตรอนปล่อย 2.5 ตัวต่อฟิชชัน จะต้องให้นิวตรอนเกิดฟิชชันกี่ตัว?
\begin{enumerate}[label=\alph*)., itemsep=0.3\baselineskip, topsep=0.3\baselineskip]
\item ก\ 
\begin{minipage}[t]{0.9\linewidth}
\raggedright 0.5 ตัว
\end{minipage}
\item ข\ 
\begin{minipage}[t]{0.9\linewidth}
\raggedright 1 ตัว
\end{minipage}
\item ค\ 
\begin{minipage}[t]{0.9\linewidth}
\raggedright 1.5 ตัว
\end{minipage}
\item ง\ 
\begin{minipage}[t]{0.9\linewidth}
\raggedright 2.5 ตัว
\end{minipage}
\end{enumerate}
\vspace{-0.3\baselineskip}
\begin{small}
\textit{คำตอบ: ข}
\end{small}

\item ถ้าพลังงานจากฟิชชัน 1 กิโลกรัมของ ²³⁵U มีค่าเท่ากับพลังงานจากการเผาไหม้กี่กิโลกรัมของถ่านหิน? (กำหนด: ถ่านหิน 3 × 10⁷ J/kg, ฟิชชัน 8 × 10¹³ J/kg)
\begin{enumerate}[label=\alph*)., itemsep=0.3\baselineskip, topsep=0.3\baselineskip]
\item ก\ 
\begin{minipage}[t]{0.9\linewidth}
\raggedright 2.7 × 10⁶ kg
\end{minipage}
\item ข\ 
\begin{minipage}[t]{0.9\linewidth}
\raggedright 2.7 × 10⁵ kg
\end{minipage}
\item ค\ 
\begin{minipage}[t]{0.9\linewidth}
\raggedright 2.7 × 10⁴ kg
\end{minipage}
\item ง\ 
\begin{minipage}[t]{0.9\linewidth}
\raggedright 2.7 × 10³ kg
\end{minipage}
\end{enumerate}
\vspace{-0.3\baselineskip}
\begin{small}
\textit{คำตอบ: ก}
\end{small}

\item วิธีการจัดการขยะนิวเคลียร์ระดับสูง (HLW) ที่นิยมใช้ในปัจจุบันคืออะไร?
\begin{enumerate}[label=\alph*)., itemsep=0.3\baselineskip, topsep=0.3\baselineskip]
\item ก\ 
\begin{minipage}[t]{0.9\linewidth}
\raggedright ฝังในทะเล
\end{minipage}
\item ข\ 
\begin{minipage}[t]{0.9\linewidth}
\raggedright เก็บในบ่อน้ำใต้ดินชั่วคราว
\end{minipage}
\item ค\ 
\begin{minipage}[t]{0.9\linewidth}
\raggedright แปลงเป็นโมลิโบเดตแก้ว (Vitrification) แล้วเก็บกักในภาชนะหลายชั้น
\end{minipage}
\item ง\ 
\begin{minipage}[t]{0.9\linewidth}
\raggedright ส่งออกไปยังประเทศอื่น
\end{minipage}
\end{enumerate}
\vspace{-0.3\baselineskip}
\begin{small}
\textit{คำตอบ: ค}
\end{small}

\item เครื่อง Tokamak ทำงานอย่างไร?
\begin{enumerate}[label=\alph*)., itemsep=0.3\baselineskip, topsep=0.3\baselineskip]
\item ก\ 
\begin{minipage}[t]{0.9\linewidth}
\raggedright ใช้แรงโน้มถ่วงอัดพลาสมา
\end{minipage}
\item ข\ 
\begin{minipage}[t]{0.9\linewidth}
\raggedright ใช้สนามแม่เหล็กกักขังพลาสมาความร้อนสูง
\end{minipage}
\item ค\ 
\begin{minipage}[t]{0.9\linewidth}
\raggedright ใช้ความดันสูงอัดนิวเคลียส
\end{minipage}
\item ง\ 
\begin{minipage}[t]{0.9\linewidth}
\raggedright ใช้เลเซอร์กระตุ้น
\end{minipage}
\end{enumerate}
\vspace{-0.3\baselineskip}
\begin{small}
\textit{คำตอบ: ข}
\end{small}

\item สาร fissile, fissile และ fertile ต่างกันอย่างไร?
\begin{enumerate}[label=\alph*)., itemsep=0.3\baselineskip, topsep=0.3\baselineskip]
\item ก\ 
\begin{minipage}[t]{0.9\linewidth}
\raggedright Fissile = เกิดฟิชชันได้ด้วยนิวตรอนช้า, Fissionable = เกิดฟิชชันได้ด้วยนิวตรอนเร็ว, Fertile = เปลี่ยนเป็น fissile ได้หลังดูดนิวตรอน
\end{minipage}
\item ข\ 
\begin{minipage}[t]{0.9\linewidth}
\raggedright ทั้งสามหมายถึงสิ่งเดียวกัน
\end{minipage}
\item ค\ 
\begin{minipage}[t]{0.9\linewidth}
\raggedright Fissile ใช้ไม่ได้กับ ²³⁵U
\end{minipage}
\item ง\ 
\begin{minipage}[t]{0.9\linewidth}
\raggedright Fertile คือสารที่เกิดฟิชชันได้ทันที
\end{minipage}
\end{enumerate}
\vspace{-0.3\baselineskip}
\begin{small}
\textit{คำตอบ: ก}
\end{small}

\item การเกิด Super-critical mass หมายถึงอะไร?
\begin{enumerate}[label=\alph*)., itemsep=0.3\baselineskip, topsep=0.3\baselineskip]
\item ก\ 
\begin{minipage}[t]{0.9\linewidth}
\raggedright มวลน้อยกว่ามวลวิกฤต
\end{minipage}
\item ข\ 
\begin{minipage}[t]{0.9\linewidth}
\raggedright มวลมากกว่ามวลวิกฤตมากจนนิวตรอนเพิ่มขึ้นแบบทวีคูณ
\end{minipage}
\item ค\ 
\begin{minipage}[t]{0.9\linewidth}
\raggedright มวลเท่ากับมวลวิกฤตพอดี
\end{minipage}
\item ง\ 
\begin{minipage}[t]{0.9\linewidth}
\raggedright ไม่มีนิวตรอนปล่อยออกมา
\end{minipage}
\end{enumerate}
\vspace{-0.3\baselineskip}
\begin{small}
\textit{คำตอบ: ข}
\end{small}

\item ปฏิกิริยา Triple-alpha คืออะไร?
\begin{enumerate}[label=\alph*)., itemsep=0.3\baselineskip, topsep=0.3\baselineskip]
\item ก\ 
\begin{minipage}[t]{0.9\linewidth}
\raggedright โปรตอน 3 ตัวรวมกัน
\end{minipage}
\item ข\ 
\begin{minipage}[t]{0.9\linewidth}
\raggedright อนุภาคแอลฟา 3 ตัวรวมกันเป็นคาร์บอน-12
\end{minipage}
\item ค\ 
\begin{minipage}[t]{0.9\linewidth}
\raggedright นิวตรอน 3 ตัวรวมกัน
\end{minipage}
\item ง\ 
\begin{minipage}[t]{0.9\linewidth}
\raggedright ฟิชชัน 3 ครั้งติดต่อกัน
\end{minipage}
\end{enumerate}
\vspace{-0.3\baselineskip}
\begin{small}
\textit{คำตอบ: ข}
\end{small}

\item วิธี Transmutation ของขยะนิวเคลียร์ทำงานอย่างไร?
\begin{enumerate}[label=\alph*)., itemsep=0.3\baselineskip, topsep=0.3\baselineskip]
\item ก\ 
\begin{minipage}[t]{0.9\linewidth}
\raggedright เผาขยะนิวเคลียร์
\end{minipage}
\item ข\ 
\begin{minipage}[t]{0.9\linewidth}
\raggedright ใช้นิวตรอนหรือโปรตอนยิงเปลี่ยนไอโซโทปกัมมันตรังสีเป็นไอโซโทปที่มีค่าครึ่งชีวิตสั้นกว่า
\end{minipage}
\item ค\ 
\begin{minipage}[t]{0.9\linewidth}
\raggedright ส่งขยะไปยังอวกาศ
\end{minipage}
\item ง\ 
\begin{minipage}[t]{0.9\linewidth}
\raggedright ใช้สนามไฟฟ้าแยกไอโซโทป
\end{minipage}
\end{enumerate}
\vspace{-0.3\baselineskip}
\begin{small}
\textit{คำตอบ: ข}
\end{small}

\item ถ้ามวลวิกฤตของ ²³⁵U ทรงกลมบริสุทธิ์ = 52 kg รัศมีวิกฤตจะประมาณเท่าใด?
\begin{enumerate}[label=\alph*)., itemsep=0.3\baselineskip, topsep=0.3\baselineskip]
\item ก\ 
\begin{minipage}[t]{0.9\linewidth}
\raggedright 3 cm
\end{minipage}
\item ข\ 
\begin{minipage}[t]{0.9\linewidth}
\raggedright 8 cm
\end{minipage}
\item ค\ 
\begin{minipage}[t]{0.9\linewidth}
\raggedright 15 cm
\end{minipage}
\item ง\ 
\begin{minipage}[t]{0.9\linewidth}
\raggedright 20 cm
\end{minipage}
\end{enumerate}
\vspace{-0.3\baselineskip}
\begin{small}
\textit{คำตอบ: ข}
\end{small}

\item เครื่องปฏิกรณ์นิวเคลียร์ที่ใช้น้ำหนัก (Heavy water) เป็นสารชะลอมีข้อดีอย่างไร?
\begin{enumerate}[label=\alph*)., itemsep=0.3\baselineskip, topsep=0.3\baselineskip]
\item ก\ 
\begin{minipage}[t]{0.9\linewidth}
\raggedright ถูกกว่าน้ำเบา
\end{minipage}
\item ข\ 
\begin{minipage}[t]{0.9\linewidth}
\raggedright ใช้ยูเรเนียมธรรมชาติได้โดยไม่ต้อง enrichment
\end{minipage}
\item ค\ 
\begin{minipage}[t]{0.9\linewidth}
\raggedright ทนความร้อนสูงกว่า
\end{minipage}
\item ง\ 
\begin{minipage}[t]{0.9\linewidth}
\raggedright มีขนาดเล็กกว่า
\end{minipage}
\end{enumerate}
\vspace{-0.3\baselineskip}
\begin{small}
\textit{คำตอบ: ข}
\end{small}

\item Gun-type bomb และ Implosion-type bomb ต่างกันอย่างไร?
\begin{enumerate}[label=\alph*)., itemsep=0.3\baselineskip, topsep=0.3\baselineskip]
\item ก\ 
\begin{minipage}[t]{0.9\linewidth}
\raggedright Gun-type ยิงอนุภาคแอลฟา Implosion-type ยิงนิวตรอน
\end{minipage}
\item ข\ 
\begin{minipage}[t]{0.9\linewidth}
\raggedright Gun-type ยิงชิ้นส่วนยูเรเนียมเข้าหากัน Implosion-type อัดระเบิดเข้าหาศูนย์กลาง
\end{minipage}
\item ค\ 
\begin{minipage}[t]{0.9\linewidth}
\raggedright ทั้งสองเหมือนกัน
\end{minipage}
\item ง\ 
\begin{minipage}[t]{0.9\linewidth}
\raggedright Gun-type ใช้พลูโทเนียม Implosion-type ใช้ยูเรเนียม
\end{minipage}
\end{enumerate}
\vspace{-0.3\baselineskip}
\begin{small}
\textit{คำตอบ: ข}
\end{small}

\item ข้อใดต่อไปนี้ไม่ใช่ข้อดีของพลังงานฟิวชันเมื่อเทียบกับฟิชชัน?
\begin{enumerate}[label=\alph*)., itemsep=0.3\baselineskip, topsep=0.3\baselineskip]
\item ก\ 
\begin{minipage}[t]{0.9\linewidth}
\raggedright มีเชื้อเพลิงอยู่มากบนโลก
\end{minipage}
\item ข\ 
\begin{minipage}[t]{0.9\linewidth}
\raggedright ไม่เกิดขยะกัมมันตรังสีระยะยาว
\end{minipage}
\item ค\ 
\begin{minipage}[t]{0.9\linewidth}
\raggedright เป็นพลังงานหมดแน่นอนในอนาคต
\end{minipage}
\item ง\ 
\begin{minipage}[t]{0.9\linewidth}
\raggedright ไม่เสี่ยงต่อการหลอมละลาย (meltdown)
\end{minipage}
\end{enumerate}
\vspace{-0.3\baselineskip}
\begin{small}
\textit{คำตอบ: ค}
\end{small}

\item ถ้าพลังงานจากฟิวชัน D-T หนึ่งครั้ง = 17.6 MeV และต้องการกำลัง 1 GW จงหาจำนวนปฏิกิริยาฟิวชันต่อวินาที
\begin{enumerate}[label=\alph*)., itemsep=0.3\baselineskip, topsep=0.3\baselineskip]
\item ก\ 
\begin{minipage}[t]{0.9\linewidth}
\raggedright 1.1 × 10²⁰ ครั้ง/s
\end{minipage}
\item ข\ 
\begin{minipage}[t]{0.9\linewidth}
\raggedright 3.5 × 10²⁰ ครั้ง/s
\end{minipage}
\item ค\ 
\begin{minipage}[t]{0.9\linewidth}
\raggedright 3.5 × 10¹⁸ ครั้ง/s
\end{minipage}
\item ง\ 
\begin{minipage}[t]{0.9\linewidth}
\raggedright 1.1 × 10¹⁸ ครั้ง/s
\end{minipage}
\end{enumerate}
\vspace{-0.3\baselineskip}
\begin{small}
\textit{คำตอบ: ข}
\end{small}

\item ไอโซโทปกัมมันตรังสีในขยะนิวเคลียร์ที่มีค่าครึ่งชีวิตสั้นมากจะมีกัมมันตภาพรังสีอย่างไร?
\begin{enumerate}[label=\alph*)., itemsep=0.3\baselineskip, topsep=0.3\baselineskip]
\item ก\ 
\begin{minipage}[t]{0.9\linewidth}
\raggedright ต่ำมาก
\end{minipage}
\item ข\ 
\begin{minipage}[t]{0.9\linewidth}
\raggedright สูงมากแต่ลดลงเร็ว
\end{minipage}
\item ค\ 
\begin{minipage}[t]{0.9\linewidth}
\raggedright คงที่ตลอดไป
\end{minipage}
\item ง\ 
\begin{minipage}[t]{0.9\linewidth}
\raggedright เป็นศูนย์
\end{minipage}
\end{enumerate}
\vspace{-0.3\baselineskip}
\begin{small}
\textit{คำตอบ: ข}
\end{small}

\item เครื่องปฏิกรณ์นิวเคลียร์แบบ PWR และ BWR ต่างกันอย่างไร?
\begin{enumerate}[label=\alph*)., itemsep=0.3\baselineskip, topsep=0.3\baselineskip]
\item ก\ 
\begin{minipage}[t]{0.9\linewidth}
\raggedright PWR ใช้น้ำเบาเป็นทั้งตัวหล่อเย็นและสารชะลอ BWR ใช้แก๊ส
\end{minipage}
\item ข\ 
\begin{minipage}[t]{0.9\linewidth}
\raggedright PWR รักษาความดันสูงไม่ให้น้ำเดือด BWR ปล่อยให้น้ำเดือดเป็นไอ
\end{minipage}
\item ค\ 
\begin{minipage}[t]{0.9\linewidth}
\raggedright PWR ดีกว่า BWR ทุกประการ
\end{minipage}
\item ง\ 
\begin{minipage}[t]{0.9\linewidth}
\raggedright ทั้งสองใช้น้ำหนักเหมือนกัน
\end{minipage}
\end{enumerate}
\vspace{-0.3\baselineskip}
\begin{small}
\textit{คำตอบ: ข}
\end{small}

\item ข้อใดต่อไปนี้เป็นเหตุผลหลักที่ทำให้ฟิชชันปล่อยพลังงานมาก?
\begin{enumerate}[label=\alph*)., itemsep=0.3\baselineskip, topsep=0.3\baselineskip]
\item ก\ 
\begin{minipage}[t]{0.9\linewidth}
\raggedright เพราะเกิดปฏิกิริยาเคมี
\end{minipage}
\item ข\ 
\begin{minipage}[t]{0.9\linewidth}
\raggedright เพราะมวลต่ำ (mass defect) ระหว่างสารตั้งต้นและผลิตภัณฑ์ถูกเปลี่ยนเป็นพลังงาน
\end{minipage}
\item ค\ 
\begin{minipage}[t]{0.9\linewidth}
\raggedright เพราะนิวเคลียสแตกตัวปล่อยอนุภาค
\end{minipage}
\item ง\ 
\begin{minipage}[t]{0.9\linewidth}
\raggedright เพราะอิเล็กตรอนถูกปล่อย
\end{minipage}
\end{enumerate}
\vspace{-0.3\baselineskip}
\begin{small}
\textit{คำตอบ: ข}
\end{small}

\end{enumerate}
\end{document}